\chapter{Grid summary}\label{ch:gridsum}
\begin{goals}
\textbf{In this chapter} you find a block model's grid from its centroids: origin, block size and count
on each axis, the order its rows are written in, sub-blocks, and how full it is.
\end{goals}

A block model's grid is rarely written down where you need it. \menu{Tools>Grid summary…} works it out
from the coordinates in the file: the numbers you type into another package's dialog to load the model,
and a check that the model is the regular grid it claims to be.

\section{Running it}

The window asks for the centroid columns \ui{X}, \ui{Y} and \ui{Z}, and optionally the block-size columns
\ui{dx}, \ui{dy} and \ui{dz}. It fills them in from columns with those names. Click \ui{Summarize}:
one pass over the view (only the filtered rows, with a filter on).

\shot[0.7\linewidth]{gridsum}{Grid summary of the sample: a 60 × 50 × 40 grid of 10 × 10 × 5~m blocks,
X varying fastest, and every cell filled.}

\section{The result}\label{sec:gridsum-result}

A table by axis:

\begin{description}
\item[\ui{origin}] the lowest block \emph{centroid} on the axis. Not the corner of the model: the corner is
half a block below, origin minus half the block size. Datamine's \col{XMORIG} and most packages' model
origins are corners.
\item[\ui{block size}] the spacing of the centroids.
\item[\ui{count}] the number of blocks along the axis, from the lowest centroid to the highest:
$(\max - \min)/\mathrm{size} + 1$.
\item[\ui{min}], \ui{max} the lowest and highest centroid.
\end{description}

and below it:

\begin{description}
\item[\ui{ordering}] which coordinate changes fastest from one row to the next and which slowest:
\code{X fastest · Y · Z slowest} is a file written X within Y within Z.
\item[\ui{parent cells}] the cells of the full grid, $n_x n_y n_z$.
\item[\ui{parent block volume}] one block's volume.
\item[\ui{rows scanned}] the rows read.
\item[\ui{fill \%}] rows scanned over parent cells. 100\% is a full grid; less, a model that stores only
some blocks (only the ore, only below topography); more than 100\%, a sub-blocked model or one with
duplicate blocks, and the window says so.
\item[\ui{X sub-blocks}] and so on, when an axis is sub-blocked: each sub-block size with the number of
them in a parent, \code{2.5 (÷4)}.
\end{description}

Every number is shown as it is, without thousands separators or rounding, and clicking it copies it, so
it pastes into another program's dialog box exactly. \ui{copy all} copies the whole summary.

\section{How it is worked out}\label{sec:gridsum-method}

For each axis, lamina collects the distinct centroid values, up to 200{,}000 of them, sorts them, and
counts the gaps between neighbours (rounded to ten significant figures, so that float noise in an
export doesn't count as a different gap).

\begin{itemize}
\item One gap size, or one that makes up more than 98\% of the gaps: that is the block size.
\item Several sizes, where the largest is a whole multiple (2 or more) of the others: a sub-blocked axis,
with the largest as the parent size and the others as sub-blocks.
\item Otherwise, the most common gap.
\end{itemize}

Block-size columns, when you give them, take precedence over the gaps: one distinct value is the block
size; several make the largest the parent and the rest its sub-blocks.

The ordering comes from counting how often each coordinate changes between consecutive rows: the one
that changes most often varies fastest.

\section{When there is no grid}

\ui{no regular grid detected} is an answer too: the coordinates don't sit on one consistent spacing,
because the data are scattered (samples, not blocks), the model is rotated (its centroids aren't on
an axis-aligned grid), or an export has gone wrong. The table then gives each axis's count of distinct
values and its range.

If an axis has more than 200{,}000 distinct values the window notes that the cap was reached; the
inference uses the values it kept.

\begin{tryit}
From the sample to another package.
\begin{enumerate}
\item Open the sample and \menu{Tools>Grid summary…}, then \ui{Summarize}.
\item Read the X origin, 620000, and the block size, 10: the model's corner in X is 619995.
\item Click 620000 to copy it. Paste it anywhere: it is the number alone.
\item Filter \code{LITO != "WASTE"} and summarise again. The grid now spans only the extent of the
mineralised blocks, and since the waste between them is filtered out, the fill is below 100\%.
\end{enumerate}
\end{tryit}
